\documentclass[tikz,border=10pt]{standalone}
\usepackage[no-math]{fontspec}
\usepackage[brazil]{babel}
\usepackage{amsmath,amssymb}
\IfFileExists{C:/Windows/Fonts/Lato_Regular.ttf}{%
  \setmainfont{Lato}[
    Path = C:/Windows/Fonts/,
    UprightFont = Lato_Regular.ttf,
    BoldFont = Lato_Bold.ttf,
    ItalicFont = Lato_Italic.ttf,
    BoldItalicFont = Lato_BoldItalic.ttf
  ]
  \setsansfont{Lato}[
    Path = C:/Windows/Fonts/,
    UprightFont = Lato_Regular.ttf,
    BoldFont = Lato_Bold.ttf,
    ItalicFont = Lato_Italic.ttf,
    BoldItalicFont = Lato_BoldItalic.ttf
  ]
}{%
  \RequirePackage[default]{lato}
}
\renewcommand{\familydefault}{\sfdefault}

\usepackage{xcolor}
\usepackage{tikz}
\usepackage{pgfplots}
\pgfplotsset{compat=1.18}
\usetikzlibrary{arrows.meta,positioning,shapes.geometric,fit,backgrounds,calc,decorations.pathreplacing,patterns}

% Paleta institucional PLANTERR / UEFS (alinhada ao estilo planterr-slides.css / 09a)
\definecolor{etNavy}{HTML}{1F4268}
\definecolor{etNavyDark}{HTML}{152E4A}
\definecolor{etGold}{HTML}{B38A3D}
\definecolor{etGoldDark}{HTML}{8C6927}
\definecolor{etTeal}{HTML}{0E7C7B}
\definecolor{etForest}{HTML}{2D6A4F}
\definecolor{etAlert}{HTML}{9E2A2B}
\definecolor{etInk}{HTML}{26313D}
\definecolor{etMuted}{HTML}{5B6573}
\definecolor{etBorder}{HTML}{D5DEE8}

\definecolor{bgNavy}{HTML}{EDF2F7}
\definecolor{bgGold}{HTML}{F8F3E9}
\definecolor{bgTeal}{HTML}{EBF6F5}
\definecolor{bgForest}{HTML}{EEF5F0}
\definecolor{bgAlert}{HTML}{FDF2F2}
\definecolor{bgCanvas}{HTML}{FFFFFF}

\tikzset{
  >=Stealth,
  etBox/.style={
    draw=#1, fill=white, rounded corners=5pt, line width=1.15pt,
    inner sep=7pt, align=center, font=\footnotesize\sffamily, text=etInk
  },
  etTint/.style={
    draw=#1, fill=#1!8!white, rounded corners=5pt, line width=1.15pt,
    inner sep=7pt, align=center, font=\footnotesize\sffamily, text=etInk
  },
  etBadge/.style={
    fill=#1, text=white, rounded corners=9pt, inner xsep=8pt, inner ysep=3pt,
    font=\scriptsize\bfseries\sffamily
  },
  etTitle/.style={font=\normalsize\bfseries\sffamily, text=etNavy, align=center},
  etSub/.style={font=\footnotesize\sffamily, text=etMuted, align=center},
  etArrow/.style={->, >=Stealth, line width=1.35pt},
  etDArrow/.style={<->, >=Stealth, line width=1.35pt}
}


\begin{document}
\begin{tikzpicture}[font=\sffamily]

% Title
\node[etTitle] (title) at (8.1, 7.65) {Business Model Canvas (9 Blocos) e Ciclo Lean Startup};

% ==================== LEFT ZONE: BUSINESS MODEL CANVAS (10.5 x 6.4 cm) ====================
\node[etSub, font=\footnotesize\bfseries\sffamily, text=etNavy] at (5.25, 6.78) {Business Model Canvas (9 Blocos)};

% Canvas block helper style
\tikzset{
  bmcCell/.style={
    draw=#1, fill=#1!8!white, line width=1pt, rounded corners=2pt
  },
  bmcHead/.style={
    font=\scriptsize\bfseries\sffamily, text=#1, align=center, inner sep=3pt
  },
  bmcBody/.style={
    font=\tiny\sffamily, text=etMuted, align=center, inner sep=2pt
  }
}

% Grid dimensions: x in [0, 10.5] (5 cols x 2.1cm), y in [0, 6.4]
% Upper rows y in [1.6, 6.4] (height 4.8cm, split at y=4.0 -> 2.4cm each)
% Bottom row y in [0, 1.6] (height 1.6cm, split at x=5.25 -> 5.25cm each)

% 1. Parcerias-Chave (col 1, full upper)
\draw[bmcCell=etNavy] (0.00, 1.60) rectangle (2.06, 6.40);
\node[bmcHead=etNavy, anchor=north, text width=1.85cm] at (1.03, 6.26) {Parcerias-Chave};
\node[bmcBody, text width=1.85cm] at (1.03, 3.85) {Rede de aliados,\\ICTs, fornecedores\\e cooperativas};

% 2. Atividades-Chave (col 2, upper)
\draw[bmcCell=etNavy] (2.11, 4.03) rectangle (4.17, 6.40);
\node[bmcHead=etNavy, anchor=north, text width=1.85cm] at (3.14, 6.26) {Atividades-Chave};
\node[bmcBody, text width=1.85cm] at (3.14, 4.85) {P\&D, produção\\e operação};

% 3. Recursos-Chave (col 2, lower)
\draw[bmcCell=etNavy] (2.11, 1.60) rectangle (4.17, 3.97);
\node[bmcHead=etNavy, anchor=north, text width=1.85cm] at (3.14, 3.83) {Recursos-Chave};
\node[bmcBody, text width=1.85cm] at (3.14, 2.45) {PI, equipe e ativos\\territoriais};

% 4. Proposta de Valor (col 3, full upper, highlighted Gold)
\draw[draw=etGold, fill=bgGold, line width=1.35pt, rounded corners=2pt] (4.22, 1.60) rectangle (6.28, 6.40);
\node[bmcHead=etGoldDark, anchor=north, text width=1.85cm] at (5.25, 6.26) {Proposta de Valor};
\node[bmcBody, text=etInk, text width=1.85cm] at (5.25, 3.85) {Dor resolvida ou\\ganho entregue\\ao cliente};

% 5. Relacionamento (col 4, upper)
\draw[bmcCell=etTeal] (6.33, 4.03) rectangle (8.39, 6.40);
\node[bmcHead=etTeal, anchor=north, text width=1.85cm] at (7.36, 6.26) {Relacionamento};
\node[bmcBody, text width=1.85cm] at (7.36, 4.85) {Aquisição, retenção\\e suporte};

% 6. Canais (col 4, lower)
\draw[bmcCell=etTeal] (6.33, 1.60) rectangle (8.39, 3.97);
\node[bmcHead=etTeal, anchor=north, text width=1.85cm] at (7.36, 3.83) {Canais};
\node[bmcBody, text width=1.85cm] at (7.36, 2.45) {Entrega física\\ou digital};

% 7. Segmentos de Clientes (col 5, full upper)
\draw[bmcCell=etTeal] (8.44, 1.60) rectangle (10.50, 6.40);
\node[bmcHead=etTeal, anchor=north, text width=1.85cm] at (9.47, 6.26) {Segmentos de Clientes};
\node[bmcBody, text width=1.85cm] at (9.47, 3.85) {Perfis e nichos\\atendidos};

% 8. Estrutura de Custos (bottom left)
\draw[bmcCell=etAlert] (0.00, 0.00) rectangle (5.22, 1.54);
\node[bmcHead=etAlert, anchor=north, text width=4.8cm] at (2.61, 1.42) {Estrutura de Custos};
\node[bmcBody, text width=4.8cm] at (2.61, 0.55) {Custos fixos, variáveis e escala operacional};

% 9. Fontes de Receita (bottom right)
\draw[bmcCell=etForest] (5.28, 0.00) rectangle (10.50, 1.54);
\node[bmcHead=etForest, anchor=north, text width=4.8cm] at (7.89, 1.42) {Fontes de Receita};
\node[bmcBody, text width=4.8cm] at (7.89, 0.55) {Precificação, recorrência e licenciamento};

% ==================== RIGHT ZONE: LEAN STARTUP LOOP ====================
\node[etSub, font=\footnotesize\bfseries\sffamily, text=etNavy] at (13.65, 6.78) {Ciclo Lean Startup};

\node[draw=etNavy, fill=bgNavy, rounded corners=6pt, line width=1.25pt,
      text width=3.55cm, minimum height=1.22cm, align=center, font=\footnotesize\sffamily, text=etInk]
  (b1) at (13.25, 5.45) {\textbf{\color{etNavy}1. Construir}\\{\scriptsize (MVP / Hipótese)}};

\node[draw=etTeal, fill=bgTeal, rounded corners=6pt, line width=1.25pt,
      text width=3.55cm, minimum height=1.22cm, align=center, font=\footnotesize\sffamily, text=etInk]
  (b2) at (13.25, 3.20) {\textbf{\color{etTeal}2. Medir}\\{\scriptsize (Métricas Acionáveis)}};

\node[draw=etGold, fill=bgGold, rounded corners=6pt, line width=1.25pt,
      text width=3.55cm, minimum height=1.22cm, align=center, font=\footnotesize\sffamily, text=etInk]
  (b3) at (13.25, 0.95) {\textbf{\color{etGoldDark}3. Aprender}\\{\scriptsize (Pivotar ou Perseverar)}};

\draw[etArrow, etNavy] (b1.south) -- (b2.north);
\draw[etArrow, etTeal] (b2.south) -- (b3.north);

% Clean external return arrow on the right edge (never behind box 2!)
\draw[etArrow, etGoldDark, line width=1.45pt, rounded corners=8pt]
  (b3.east) -- (16.15, 0.95) -- (16.15, 5.45) -- (b1.east);

\node[fill=white, text=etGoldDark, font=\scriptsize\bfseries\sffamily, rotate=90, inner sep=3pt]
  at (16.15, 3.20) {Iteração Contínua};

% Subtle link from Canvas to Lean Loop
\draw[etArrow, etMuted, dashed, line width=1.15pt] (10.65, 3.20) -- (b2.west);

\end{tikzpicture}
\end{document}
